\chapter{Manjul Bhargava}\label{ch2014manjulbhargava:manjul-bhargava-fields-medal-2014-a-comprehensive-report-on-his-work-and-impact}

\begin{center}\rule{0.5\linewidth}{0.5pt}\end{center}

\textbf{Prepared:} July 2026 \textbf{Subject:} Manjul Bhargava (b. August 8, 1974, Hamilton, Ontario, Canada) \textbf{Position:} Robert C. Gunning and R. Brandon Fradd Professor of Mathematics, Princeton University \textbf{Doctoral Advisor:} Andrew Wiles \textbf{Fields Medal Citation:} \emph{``For developing powerful new methods in the geometry of numbers, which he applied to count rings of small rank and to bound the average rank of elliptic curves.''}

\begin{center}\rule{0.5\linewidth}{0.5pt}\end{center}

\begin{prereqs} Elementary number theory (quadratic forms, class numbers), abstract algebra (rings of integers, group actions), basic algebraic geometry. \end{prereqs}

\begin{objectives} Explain higher composition laws and the Bhargava boxes \cite{bhargava2015}; describe the average size of 2-Selmer groups; understand the connection between number fields and the geometry of numbers. \end{objectives}

\section{Problem History}

In 1801, Gauss discovered his celebrated law of composition for binary quadratic forms in the \emph{Disquisitiones Arithmeticae}. A binary quadratic form $Q(x,y) = ax^2 + bxy + cy^2$ can be ``composed'' with another to produce a third form of the same discriminant, and this operation gives the set of equivalence classes of forms a group structure isomorphic to the ideal class group of the quadratic field $\mathbb{Q}(\sqrt{D})$. Gauss's proof required 20 pages of dense calculations, and for two centuries this was considered the definitive treatment.

Mathematicians long wondered whether analogous composition laws existed for forms of higher degree. Binary cubic forms, ternary quadratic forms, and higher-degree forms were studied, but no general composition laws were found. The prevailing view became that Gauss composition was a unique phenomenon tied to the special structure of binary quadratic forms.

The problem of counting number fields of small degree was another classical question. The Minkowski bound implies that the number $N_n^{\pm}(X)$ of degree-$n$ number fields with absolute discriminant at most $X$ grows at most like $X^{n/2}$. Precise asymptotics were known for $n=2$ (quadratic fields, by Gauss) and $n=3$ (cubic fields, by Davenport and Heilbronn in the 1960s--70s, who proved $N_3^{\pm}(X) \sim X/(3\zeta(3))$ and $X/(4\zeta(3))$ for totally real and complex cases respectively). For $n=4$ and $n=5$, the problem had been open for over 50 years.

The distribution of ranks of elliptic curves over $\mathbb{Q}$ was perhaps the deepest mystery. The Birch and Swinnerton-Dyer conjecture, one of the seven Clay Millennium Prize Problems, predicts that the rank of an elliptic curve $E$ equals the order of vanishing of its $L$-function at $s=1$. The Goldfeld--Katz--Sarnak ``minimalist conjecture'' predicts that 50\% of elliptic curves have rank 0, 50\% have rank 1, and essentially none have rank $\ge 2$, giving an average rank of $1/2$. Before Bhargava's work, it was not even known that the average rank is finite.

The connection between these problems became clear through Bhargava's higher composition laws, which provided the missing parametrizations needed to count quartic and quintic rings and fields. His work on Selmer groups then translated this counting technology into bounds on the average rank of elliptic curves.

The fundamental insight is that counting algebraic objects---number fields, rings, elliptic curves---requires parametrizing them by orbits of group actions on spaces of forms. Once such a parametrization is established, one can apply the geometry of numbers (Minkowski's theorem, Siegel's integral formula) to count orbits and hence count the algebraic objects. Before Bhargava, this strategy was known only for quadratic and cubic forms; Bhargava extended it to quartic and quintic forms, opening the door to previously inaccessible counting problems.

\section{How Was It Solved}

Bhargava's breakthrough began with a geometric reinterpretation of Gauss composition. While playing with a $2\times 2\times 2$ Rubik's cube in the summer of 1998, he realized that labeling the eight corners of a cube with integers and slicing it along each of its three axes produces three pairs of $2\times 2$ matrices. From each pair $(M,N)$, one constructs a binary quadratic form $Q_i(x,y) = -\det(M_ix + N_iy)$. The three forms $Q_1, Q_2, Q_3$ have the same discriminant, and $Q_1 \circ Q_2 = Q_3$ in the sense of Gauss composition. This gave a remarkably simple proof of Gauss's 20-page composition law.

This insight revealed that Gauss composition was just the first instance of a much more general family. By considering cubes and hypercubes of varying dimensions, Bhargava discovered 14 new higher composition laws, published in a celebrated series of four Annals of Mathematics papers. The laws included cubic compositions from $3\times 3\times 3$ cubes and higher-degree compositions from larger combinatorial structures.

For quartic rings---commutative rings free of rank 4 as $\mathbb{Z}$-modules---Bhargava proved that they are parametrized by $\mathrm{GL}_2(\mathbb{Z})\times\mathrm{SL}_3(\mathbb{Z})$-orbits of pairs of ternary quadratic forms $(A,B)$. This generalization of the Delone--Faddeev correspondence (which parametrizes cubic rings by binary cubic forms) was the crucial ingredient for counting quartic fields. The discriminant of the quartic ring is a polynomial in the coefficients of $A$ and $B$.

For quintic rings, Bhargava established a bijection between $\mathrm{GL}_4(\mathbb{Z})\times\mathrm{SL}_5(\mathbb{Z})$-orbits of quadruples of alternating $5\times 5$ bilinear forms and isomorphism classes of quintic rings. This was used to prove the first asymptotic formula for quintic fields: $N_5^{\pm}(X) \sim C_5^{\pm} X$, where $C_5^{\pm}$ are explicit constants approximately equal to $0.0585$ (totally real) and $0.9363$ (totally complex).

The counting of quintic fields proceeds in two steps: first, count quintic rings (commutative rings free of rank 5 as $\mathbb{Z}$-modules) with bounded discriminant, and second, count the number of fields with a given ring of integers. The first step uses the parametrization by quadruples of alternating forms and the geometry of numbers, while the second step uses Minkowski's bound on the number of rings with bounded discriminant.

The asymptotic formula $N_5^{\pm}(X) \sim C_5^{\pm} X$ shows that the number of quintic fields with discriminant at most $X$ grows linearly in $X$. This is the same growth rate as for cubic fields (Davenport--Heilbronn), but the constant $C_5^{\pm}$ is much smaller because the parametrization of quintic rings is more complex. The explicit computation of $C_5^{\pm}$ requires evaluating integrals over fundamental domains of the action of $\mathrm{GL}_4(\mathbb{Z})\times\mathrm{SL}_5(\mathbb{Z})$ on spaces of alternating forms, which is a difficult problem in the geometry of numbers.

The work on elliptic curves was carried out with Arul Shankar. The key idea is that the $p$-Selmer group $\mathrm{Sel}_p(E)$ sits in the exact sequence
\[
0 \to E(\mathbb{Q})/pE(\mathbb{Q}) \to \mathrm{Sel}_p(E) \to \mathrm{Sha}(E)[p] \to 0,
\]
so $\operatorname{rank}(E) \le \dim_{\mathbb{F}_p}\mathrm{Sel}_p(E)$. Bhargava and Shankar proved that elements of $\mathrm{Sel}_p(E)$ correspond to orbits of group actions on spaces of forms of degree $p$: binary quartic forms for $\mathrm{Sel}_2$, ternary cubics for $\mathrm{Sel}_3$, and higher-dimensional spaces for $\mathrm{Sel}_4$ and $\mathrm{Sel}_5$. They counted these orbits using geometry-of-numbers methods to obtain:
\[
\text{average size of } \mathrm{Sel}_2(E) = 3, \quad \mathrm{Sel}_3(E) = 4, \quad \mathrm{Sel}_4(E) = 7, \quad \mathrm{Sel}_5(E) = 6.
\]

The geometry-of-numbers counting proceeds by showing that the set of orbits satisfies a uniform bound on the ``essential dimension'' of the parameter space. Using Siegel's integral formula and Minkowski's theorem on lattice points in convex bodies, one shows that the average size of the Selmer group is finite. The key technical innovation is a reduction to counting lattice points in fundamental domains of the action of $\mathrm{GL}_n(\mathbb{Z})$ on spaces of forms.

The Selmer group $\mathrm{Sel}_p(E)$ sits inside the larger Shafarevich--Tate group $\mathrm{Sha}(E)$, whose elements represent obstruction to the Hasse principle. Understanding the average size of $\mathrm{Sha}(E)$ is one of the deepest open problems in arithmetic geometry. The Selmer group bounds provide upper bounds on the average rank, but they do not directly bound the average size of $\mathrm{Sha}(E)$. This distinction is important because the Selmer group can be larger than the Mordell--Weil group due to the presence of nontrivial elements in $\mathrm{Sha}(E)[p]$.

The averaging over all elliptic curves is performed using the parametrization of elliptic curves by pairs $(A,B)$ of binary quartic forms (for the 2-Selmer group). The space of such pairs is acted on by $\mathrm{GL}_2(\mathbb{Z})$, and the average is taken over a fundamental domain for this action. The key estimate is that the number of orbits with bounded invariants grows polynomially, which allows the application of Siegel's weighted integral formula to compute the average.

The computation of the average Selmer group size requires careful analysis of the geometry of the fundamental domain. For the 2-Selmer group, the relevant space is the space of pairs of binary quartic forms, and the fundamental domain for the $\mathrm{GL}_2(\mathbb{Z})$ action has a complicated structure. The key technical innovation is to decompose this space into simpler pieces, count the orbits in each piece separately, and then reassemble the counts to obtain the total average. This decomposition relies on the theory of integral quadratic forms and the geometry of numbers.

The results on average Selmer group sizes have implications for the study of arithmetic statistics more generally. The methods developed by Bhargava and Shankar can be applied to other arithmetic objects, such as class groups of quadratic fields, distributions of $L$-values, and the statistics of Galois representations. The orbit-counting framework provides a unified approach to these diverse problems.

\section{Impact of This Discovery}

From the Selmer group bounds, Bhargava and Shankar deduced that the average rank of all elliptic curves over $\mathbb{Q}$ is less than $0.885$, at least $83.75\%$ of elliptic curves have rank 0 or 1, and at least $20.62\%$ have rank 0. Combined with results of Skinner--Urban and Kolyvagin, this implies that a positive proportion of elliptic curves satisfy the Birch and Swinnerton-Dyer conjecture.

The creation of \textbf{arithmetic statistics} as a subfield is perhaps Bhargava's most enduring legacy. Before his work, the distribution of number fields, class groups, and Selmer groups was studied primarily through heuristics and conjectures. Bhargava's methods made it possible to prove precise theorems about these distributions.

His higher composition laws have found unexpected applications in theoretical physics. The ``charge cube'' of STU black holes in string theory has the same structure as Bhargava's cube, and the higher composition laws describe the U-duality orbits of black hole solutions, revealing a deep link between number theory and quantum gravity.

The Bhargava factorial, introduced as an undergraduate, generalized the classical factorial function to arbitrary subsets of the integers and has found applications in the theory of integer-valued polynomials, fixed divisors of polynomials, and ring theory. For example, if $S = 2\mathbb{Z}$ (even integers), then $k!_S = 2^k k!$.

The 290-Theorem, proved with Jonathan Hanke, states that a positive-definite quadratic form with integer coefficients is universal if and only if it represents every integer up to 290. This dramatically extended the Conway--Schneeberger 15-theorem and has become a central result in the arithmetic theory of quadratic forms.

The connection between Bhargava's work and the theory of elliptic curves has implications for the study of special values of $L$-functions. The Birch and Swinnerton-Dyer conjecture predicts that the order of vanishing of $L(E,s)$ at $s=1$ equals the rank of $E$, and the leading coefficient is given by a formula involving the Regulator, the size of the Shafarevich--Tate group, and other arithmetic invariants. Bhargava's bound on the average rank provides evidence for this conjecture by showing that most elliptic curves have small rank.

The \textbf{Goldfeld--Katz--Sarnak conjecture} predicts that the rank of a typical elliptic curve is 0 or 1, with equal proportions. Bhargava's result that the average rank is less than 0.885 provides strong evidence for this conjecture, but proving it would require showing that the average rank is exactly 1/2. This remains a major open problem in arithmetic geometry.

The \textbf{ Cohen--Lenstra heuristics} predict the distribution of class groups of quadratic fields. Bhargava's work has confirmed these heuristics in several new cases, providing the first rigorous evidence for these conjectures beyond the known cases. The techniques involve counting rings of small rank and relating their class groups to the class groups of number fields. The connection between Bhargava's orbit-counting methods and the Cohen--Lenstra heuristics is an active area of research, with potential applications to the study of class groups of higher-degree number fields.

\section{Modern Views and Improvements}

Bhargava's parametrizations currently work up to degree 5. Extending these counting results to degree 6 and beyond remains a major open problem. The obstacle is that the relevant parametrizations become increasingly complex, and the geometry-of-numbers estimates are more difficult in higher dimensions.

For $p \ge 7$, the parametrizations of Selmer groups are more complex and the counting problems higher-dimensional. Determining the average size of $\mathrm{Sel}_p(E)$ for all $p$ would give even sharper bounds on the average rank and potentially prove the minimalist conjecture.

The connection between Bhargava's orbit-counting methods and the theory of algebraic stacks and moduli spaces has been recognized. His counting can be reinterpreted as counting rational points on certain moduli stacks, connecting arithmetic statistics to modern algebraic geometry.

The theory of algebraic stacks provides a framework for counting algebraic objects by counting rational points on moduli stacks. The moduli stack of elliptic curves, for example, is an algebraic stack whose rational points correspond to isomorphism classes of elliptic curves. Bhargava's counting methods can be reinterpreted as counting rational points on these moduli stacks, which provides a connection between arithmetic statistics and modern algebraic geometry.

Bhargava's work on the Hasse principle for plane cubics and hyperelliptic curves has opened the statistical study of local-to-global principles. Understanding the proportion of varieties that satisfy the Hasse principle and the Brauer--Manin obstruction remains an active area.

The techniques developed by Bhargava---counting integral orbits via the geometry of numbers, parametrizing algebraic structures by forms, and analyzing Selmer groups via representation theory---have become standard tools in arithmetic geometry, used by researchers worldwide.

The \textbf{Bhargava factorial} generalizes the classical factorial to arbitrary subsets $S$ of the nonnegative integers. For a set $S$, the Bhargava factorial $k!_S$ is defined as the product of all elements of $S$ that are less than or equal to $k$. For example, if $S = \{0, 2, 4, 6, \ldots\}$ (even nonnegative integers), then $k!_S = \prod_{j=0, 2j \le k} 2j$. The Bhargava factorial has applications to the theory of integer-valued polynomials: the polynomial $\binom{x}{n} = x(x-1)\cdots(x-n+1)/n!$ can be generalized to $\binom{x}{n}_S = \prod_{s \in S, s < n} (x-s)/(n-s)$, and the Bhargava factorial appears in the denominator of these generalized binomial coefficients.

The \textbf{290-Theorem} states that a positive-definite quadratic form $Q(x_1, \ldots, x_n) = \sum_{i \le j} a_{ij} x_i x_j$ with integer coefficients is universal (represents all positive integers) if and only if it represents all integers up to 290. The proof proceeds in two steps: first, show that the set of universal forms is closed under a natural equivalence relation, and second, check the condition for a finite number of equivalence classes. The key insight is that the genus of a quadratic form determines its local behavior, and the spinor genus determines its global behavior. The 290 bound comes from the largest integer that can distinguish between forms in different spinor genera.

Recent work has focused on the distribution of class groups of quadratic fields, where Bhargava's methods give sharp bounds on the average size of the $p$-torsion. The connection between class group statistics and Cohen--Lenstra heuristics is an active area of research, with Bhargava's methods providing the first rigorous results in many settings.

The Cohen--Lenstra heuristics predict that the distribution of $p$-parts of class groups of quadratic fields follows specific probability distributions related to random matrix theory. Bhargava's work has confirmed these heuristics in several new cases, providing the first rigorous evidence for these conjectures beyond the known cases. The techniques involve counting rings of small rank and relating their class groups to the class groups of number fields.

Another active direction is the statistical study of Galois representations. The average behavior of the Selmer groups of Galois representations associated to modular forms is being investigated using Bhargava's orbit-counting methods. This connects the arithmetic statistics program to the Langlands program and the study of automorphic forms. The development of new parametrizations for higher-rank Galois representations is a key technical challenge that, if overcome, would lead to new results on the average behavior of modular forms and their associated L-functions.

The connection between Bhargava's work and the theory of elliptic curves has implications for the BSD conjecture. The average rank bound provides evidence that most elliptic curves have small rank, which is consistent with the BSD prediction. However, proving the BSD conjecture for a positive proportion of elliptic curves remains a major open problem that requires techniques beyond the current orbit-counting framework.

The orbit-counting methods developed by Bhargava have been applied to other problems in arithmetic statistics. The distribution of class groups of number fields, the statistics of Galois representations, and the study of local-to-global principles are all areas where Bhargava's methods have led to new results. The connection between orbit-counting and the geometry of numbers continues to produce new insights into the arithmetic of algebraic objects.

The parametrization of quartic and quintic rings by Bhargava has led to new results on the distribution of these rings. The number of quartic rings with bounded discriminant grows like $X$, and the constants in the asymptotic formula have been computed explicitly. These results provide information about the arithmetic of number fields of degree 4 and 5, which had been inaccessible before Bhargava's work.

The Bhargava factorial continues to find new applications. In the theory of integer-valued polynomials, the Bhargava factorial determines the denominators of generalized binomial coefficients, which are the fundamental building blocks of integer-valued polynomial rings. The connection between the Bhargava factorial and the theory of $p$-adic analysis has led to new results on the structure of rings of integer-valued polynomials and their applications to algebra and number theory.

The connection between Bhargava's work and the Langlands program is an active area of research. The Selmer group counts that Bhargava and Shankar computed are related to the values of $L$-functions of elliptic curves, which are central objects in the Langlands program. Understanding the distribution of these $L$-values and their relationship to the arithmetic of elliptic curves is a major open problem that connects Bhargava's orbit-counting methods to the deepest questions in number theory.

The study of the distribution of class groups of number fields using Bhargava's methods has led to new evidence for the Cohen-Lenstra heuristics. The Cohen-Lenstra heuristics predict the distribution of $p$-parts of class groups of quadratic fields, and Bhargava's work has confirmed these predictions in several new cases. The connection between class group statistics and the geometry of number fields is an active area of research that builds on Bhargava's foundational work.

The average size of the $p$-torsion of the class group of quadratic fields has been studied using Bhargava's orbit-counting methods. For $p=2$, the average size of the 2-torsion is $3/2$, which is consistent with the Cohen-Lenstra predictions. For $p=3$, the average size of the 3-torsion is $4/3$, which is also consistent with the predictions. These results provide the first rigorous evidence for the Cohen-Lenstra heuristics beyond the known cases and have led to new conjectures about the distribution of class groups of higher-degree number fields.

The connection between Bhargava's work and the theory of integral points on algebraic varieties is another active direction. The Mordell-Lang conjecture predicts the structure of subgroups of the group of rational points on a variety of general type, and Bhargava's methods have been applied to special cases of this conjecture. The study of integral points on algebraic varieties using orbit-counting methods is a growing area that connects arithmetic geometry to the geometry of numbers.

The development of new parametrizations for counting algebraic objects remains a central challenge in Bhargava's program. The parametrization of quartic and quintic rings by orbits of group actions on spaces of forms has been extended to other algebraic objects, such as étale algebras and Galois representations. The development of new parametrizations for higher-rank Galois representations would have implications for the study of the Langlands program and for the distribution of modular forms.

The study of the distribution of $L$-values of elliptic curves using Bhargava's orbit-counting methods is an active area. The central value of the $L$-function of an elliptic curve is related to the size of its Selmer group, and the average size of the Selmer group provides information about the distribution of $L$-values. The development of new techniques for computing the average size of the Selmer group for larger primes $p$ would have implications for the study of the BSD conjecture and for the distribution of ranks of elliptic curves.

\begin{center}\rule{0.5\linewidth}{0.5pt}\end{center}

\section*{Summary}
\addcontentsline{toc}{section}{Summary}

Bhargava discovered 14 higher composition laws generalizing Gauss composition, developed new methods for counting number fields of small degree, and proved that the average rank of elliptic curves over $\mathbb{Q}$ is bounded, creating the field of arithmetic statistics.
